\documentclass[a4paper,12pt]{article}
\usepackage[utf8]{inputenc}
\usepackage[T1]{fontenc}
\usepackage[ngerman]{babel}
\usepackage{amsmath}
\usepackage{amssymb}
\usepackage{enumitem}
\usepackage{geometry}
\geometry{a4paper, margin=2.5cm}
\usepackage{parskip}

\title{Einsendeaufgaben -- Aufgabe 3.1\\[0.5em]
\large Modul 61112: Lineare Algebra}
\author{}
\date{}

\begin{document}
\maketitle

\section*{Aufgabe 3.1}

Es muss $\gamma^{(1)}(x) = \beta^{(1)}(f(x))$ gelten, wenn $\gamma(x,y) = \beta(f(x), y)$ für ein $f: V \to V$ Voraussetzung ist.

Da jeder Vektorraum eine Basis $B$ besitzt, wählen wir die geordnete Basis $B = (v_1, \ldots, v_n)$ aus. Da $\beta$ regulär ist, ist nach Satz 3.1.38 (ii) $\beta^{(1)}$ ein Isomorphismus. Es existiert also ein eindeutiges $x_i \in V$ mit $\beta^{(1)}(x_i) = \gamma^{(1)}(v_i)$ für $1 \leq i \leq \dim(V)$.

Wir definieren $f: V \to V$ mit
\[
  f(v_i) = x_i.
\]

Es gilt dann
\[
  \beta^{(1)}(f(v_i)) = \beta^{(1)}(x_i) = \gamma^{(1)}(v_i).
\]

Es existiert kein anderes $f': V \to V$, da $x_i$ eindeutig bestimmt ist und $f(v_i) = x_i$ gelten muss.

\end{document}
